\documentclass[UTF8,12pt,a4paper,fontset=fandol]{ctexart}

\usepackage[a4paper,top=23mm,bottom=24mm,left=25mm,right=25mm]{geometry}
\usepackage{amsmath,amssymb,bm,mathtools}
\usepackage{booktabs,longtable,array,enumitem}
\usepackage{xcolor,hyperref,fancyhdr,microtype}
\usepackage{listings}

\hypersetup{
  colorlinks=true,
  linkcolor=blue!55!black,
  urlcolor=blue!55!black
}

\setlength{\parindent}{2em}
\setlength{\parskip}{0.30em}
\setlist[itemize]{leftmargin=2.2em,itemsep=0.18em,topsep=0.30em}
\setlist[enumerate]{leftmargin=2.4em,itemsep=0.22em,topsep=0.30em}
\allowdisplaybreaks[2]
\numberwithin{equation}{section}

\pagestyle{fancy}
\setlength{\headheight}{15pt}
\fancyhf{}
\fancyhead[C]{第一题：二维对流方程完整详细解答}
\fancyfoot[C]{\thepage}
\renewcommand{\headrulewidth}{0.4pt}
\renewcommand{\footrulewidth}{0pt}

\newcolumntype{L}[1]{>{\raggedright\arraybackslash}p{#1}}
\newcommand{\Co}{\mathrm{Co}}
\newcommand{\Cmax}{\mathrm{Co}_{\max}}
\newcommand{\dd}{\,\mathrm d}
\newcommand{\Sf}{\bm S}
\newcommand{\xf}{\bm x}

\lstdefinestyle{algorithm}{
  basicstyle=\ttfamily\small,
  frame=single,
  numbers=left,
  numberstyle=\tiny,
  stepnumber=1,
  numbersep=7pt,
  showstringspaces=false,
  breaklines=true,
  columns=fullflexible
}

\title{\bfseries 第一题：二维对流方程\\从连续模型到有限体积求解器与结果诊断的完整详细解答}
\author{}
\date{}

\begin{document}

\maketitle
\vspace{-2.2em}

\begin{abstract}
本文对二维标量对流方程
\[
  \frac{\partial\phi}{\partial t}+\nabla\cdot(\bm u\phi)=0
\]
给出从零开始的系统解答。内容依次覆盖连续方程的物理意义、特征线与两个验证算例、初始条件和边界条件、有限体积半离散、一阶与二阶迎风重构、显式时间推进、Courant数、数值耗散与色散、网格分辨率 $N$、三角形与四边形网格、误差范数、收敛阶、后处理图像和异常诊断，并给出可落实为OpenFOAM或自编程序的算法流程。文中未虚构尚未计算的数值数据；所有结果部分均明确区分理论精确结果、离散格式的预期趋势和实际程序运行后应填写的量。
\end{abstract}

\tableofcontents
\newpage

\section{题目与解答目标}

题目要求求解

\begin{equation}
  \frac{\partial\phi}{\partial t}
  +\nabla\cdot(\bm u\phi)=0,
  \label{eq:governing}
\end{equation}

基于有限体积法写出半离散形式，并采用显式空间离散与显式时间推进构建数值求解器。随后用两个二维算例验证：

\begin{enumerate}
  \item 恒定速度下的正弦波周期平移，计算 $t=1$ 时的 $L_1$ 误差和网格收敛阶；
  \item 开槽圆盘、圆锥和光滑凸峰在旋转速度场中的刚体旋转，观察 $t=2\pi$ 时的轮廓保持能力。
\end{enumerate}

完整解答需要建立以下闭环：

\begin{equation}
\boxed{
\text{连续方程}
\rightarrow
\text{有限体积守恒}
\rightarrow
\text{数值通量}
\rightarrow
\text{显式时间推进}
\rightarrow
\text{误差与结果诊断}
}.
\end{equation}

\section{连续对流方程的数学与物理意义}

\subsection{标量、速度和对流通量}

$\phi(\xf,t)$ 是被输运的标量，可以表示染料浓度、污染物浓度、温度标记或其他被动标量。$\bm u(\xf,t)$ 是给定速度场。

乘积

\begin{equation}
  \bm u\phi
\end{equation}

是标量的对流通量密度。若某个有向面的单位法向量为 $\bm n$，则

\begin{equation}
  (\bm u\phi)\cdot\bm n
\end{equation}

表示单位面积、单位时间穿过该面的有向标量通量。

\subsection{守恒形式与输运形式}

利用乘积法则，

\begin{equation}
  \nabla\cdot(\bm u\phi)
  =
  \bm u\cdot\nabla\phi
  +\phi\nabla\cdot\bm u.
\end{equation}

因此式\eqref{eq:governing}可以写成

\begin{equation}
  \frac{\partial\phi}{\partial t}
  +\bm u\cdot\nabla\phi
  +\phi\nabla\cdot\bm u=0.
\end{equation}

当速度场不可压缩，

\begin{equation}
  \nabla\cdot\bm u=0,
\end{equation}

方程化为

\begin{equation}
  \frac{\partial\phi}{\partial t}
  +\bm u\cdot\nabla\phi=0.
  \label{eq:advective_form}
\end{equation}

式\eqref{eq:advective_form}表明，标量随速度场移动，但沿流体轨迹保持不变。

\subsection{特征线解释}

令粒子轨迹满足

\begin{equation}
  \frac{\dd\xf(t)}{\dd t}
  =\bm u(\xf(t),t).
  \label{eq:characteristic}
\end{equation}

沿轨迹对 $\phi$ 求全导数：

\begin{align}
  \frac{\dd}{\dd t}\phi(\xf(t),t)
  &=
  \frac{\partial\phi}{\partial t}
  +\frac{\dd\xf}{\dd t}\cdot\nabla\phi \notag\\
  &=
  \frac{\partial\phi}{\partial t}
  +\bm u\cdot\nabla\phi.
\end{align}

由式\eqref{eq:advective_form}，

\begin{equation}
  \frac{\dd}{\dd t}\phi(\xf(t),t)=0.
\end{equation}

所以连续方程本身只搬运形状，不会使形状变模糊、峰值降低、槽口填平或产生振荡。第一题没有物理扩散项，计算中出现的平滑和展宽主要来自离散格式的人工扩散。

\subsection{积分守恒律}

对固定控制体 $\Omega$ 积分：

\begin{equation}
  \frac{\dd}{\dd t}
  \int_\Omega\phi\dd V
  +
  \int_{\partial\Omega}
  \phi\bm u\cdot\bm n\dd S
  =0.
  \label{eq:integral_conservation}
\end{equation}

式\eqref{eq:integral_conservation}表示：控制体内标量总量的变化等于通过边界流出的净通量的相反数。这正是有限体积法的出发点。

\section{两个验证算例的连续精确结构}

\subsection{算例一：正弦波平移}

初始条件为

\begin{equation}
  \phi_0(x,y)=\sin 2\pi(x+y),
\end{equation}

速度场为

\begin{equation}
  \bm u=(1,1).
\end{equation}

特征线满足

\begin{equation}
  x(t)=x_0+t,
  \qquad
  y(t)=y_0+t.
\end{equation}

反解初始位置：

\begin{equation}
  x_0=x-t,
  \qquad
  y_0=y-t.
\end{equation}

沿特征线 $\phi$ 不变，因此精确解为

\begin{align}
  \phi_{\mathrm{exact}}(x,y,t)
  &=\phi_0(x-t,y-t)\notag\\
  &=\sin 2\pi(x+y-2t).
  \label{eq:sine_exact}
\end{align}

计算区域为 $[0,1]^2$，两个方向均为周期边界。$t=1$ 时，

\begin{equation}
  \phi_{\mathrm{exact}}(x,y,1)
  =
  \sin 2\pi(x+y-2)
  =
  \sin 2\pi(x+y),
\end{equation}

所以精确解完全回到初始场。该算例的价值在于：解析解光滑、已知且容易计算，适合定量评价误差与收敛阶。

\subsection{算例二：复杂轮廓刚体旋转}

速度场为

\begin{equation}
  \bm u=(0.5-y,\;x-0.5).
  \label{eq:rotation_velocity}
\end{equation}

令

\begin{equation}
  X=x-0.5,
  \qquad
  Y=y-0.5.
\end{equation}

则特征方程为

\begin{equation}
  \frac{\dd X}{\dd t}=-Y,
  \qquad
  \frac{\dd Y}{\dd t}=X.
\end{equation}

解为

\begin{equation}
  \begin{bmatrix}X(t)\\Y(t)\end{bmatrix}
  =
  \begin{bmatrix}
  \cos t&-\sin t\\
  \sin t&\cos t
  \end{bmatrix}
  \begin{bmatrix}X_0\\Y_0\end{bmatrix}.
\end{equation}

因此它是绕 $(0.5,0.5)$ 的逆时针刚体旋转，角速度为 $1$，周期为

\begin{equation}
  T=2\pi.
\end{equation}

若要在任意时刻计算精确场，应将当前位置反向旋转到初始位置：

\begin{equation}
  \begin{bmatrix}X_0\\Y_0\end{bmatrix}
  =
  \begin{bmatrix}
  \cos t&\sin t\\
  -\sin t&\cos t
  \end{bmatrix}
  \begin{bmatrix}X\\Y\end{bmatrix},
\end{equation}

然后取

\begin{equation}
  \phi_{\mathrm{exact}}(x,y,t)
  =
  \phi_0(0.5+X_0,\;0.5+Y_0).
\end{equation}

$t=2\pi$ 时旋转矩阵为单位矩阵，所以精确场应完全恢复初始形状。

\subsubsection{三个初始轮廓}

设三个轮廓半径均为

\begin{equation}
  r_0=0.15.
\end{equation}

开槽圆盘中心为 $(0.5,0.75)$，定义

\begin{equation}
  r_D=\sqrt{(x-0.5)^2+(y-0.75)^2}.
\end{equation}

取

\begin{equation}
  \phi_D(x,y)=
  \begin{cases}
  1,
  &r_D\leq0.15
  \ \text{且}\
  \left(|x-0.5|\geq0.025\ \text{或}\ y\geq0.85\right),\\
  0,&\text{其他}.
  \end{cases}
\end{equation}

圆锥中心为 $(0.5,0.25)$，定义

\begin{equation}
  r_C=\sqrt{(x-0.5)^2+(y-0.25)^2},
\end{equation}

取

\begin{equation}
  \phi_C(x,y)=
  \begin{cases}
  1-r_C/0.15,&r_C\leq0.15,\\
  0,&r_C>0.15.
  \end{cases}
\end{equation}

光滑凸峰中心为 $(0.25,0.5)$，定义

\begin{equation}
  r_H=\sqrt{(x-0.25)^2+(y-0.5)^2},
\end{equation}

取

\begin{equation}
  \phi_H(x,y)=
  \begin{cases}
  \dfrac14\left[1+\cos\left(\pi r_H/0.15\right)\right],
  &r_H\leq0.15,\\[2mm]
  0,&r_H>0.15.
  \end{cases}
\end{equation}

完整初值为

\begin{equation}
  \phi_0=\phi_D+\phi_C+\phi_H.
\end{equation}

开槽圆盘检验间断和窄槽保持能力，圆锥检验尖峰与斜面保持能力，光滑凸峰检验光滑结构上的幅值和相位误差。

\section{初始条件与边界条件的作用}

\subsection{初始条件}

初始条件

\begin{equation}
  \phi(\xf,0)=\phi_0(\xf)
\end{equation}

给出计算开始时每个位置的标量。控制方程决定“怎样运动”，初始条件决定“被运动的形状是什么”。同一个求解器只需更换初始场，就可以完成正弦波和复杂轮廓两个算例。

\subsection{双曲问题的入流与出流边界}

对于边界外法向量 $\bm n$：

\begin{itemize}
  \item $\bm u\cdot\bm n<0$ 时，信息从外部进入，是入流边界，必须规定外部标量；
  \item $\bm u\cdot\bm n>0$ 时，信息从内部流出，是出流边界，面值由内部上游状态决定。
\end{itemize}

这与扩散方程不同。纯对流方程的信息具有方向性，一般只在入流边界强制给值。

\subsection{正弦波的周期边界}

正弦波需要

\begin{equation}
  \phi(0,y,t)=\phi(1,y,t),
  \qquad
  \phi(x,0,t)=\phi(x,1,t).
\end{equation}

数值实现中必须把周期面正确配对，使一侧的出流通量成为另一侧的入流通量。周期面通量不匹配会导致边界附近误差和总量漂移。

\subsection{旋转算例的边界说明}

原题未明确给出第二算例的边界条件。三个轮廓在旋转过程中均位于以 $(0.5,0.5)$ 为中心、半径不超过约 $0.4$ 的区域内，因此不会接触外边界。可以对入流段规定背景值 $\phi=0$，出流段采用内部上游值，并在报告中明确说明。由于有效标量远离边界，该选择在一周旋转内不应影响主体轮廓。

\section{有限体积半离散的详细推导}

\subsection{控制体与单元平均}

将区域划分为不重叠控制体 $\Omega_P$。定义

\begin{equation}
  V_P=|\Omega_P|,
\end{equation}

以及单元平均

\begin{equation}
  \phi_P(t)
  =
  \frac1{V_P}
  \int_{\Omega_P}\phi(\xf,t)\dd V.
  \label{eq:cell_average}
\end{equation}

二维问题中 $V_P$ 是单元面积；OpenFOAM通常使用单位厚度或薄层三维网格，因此程序仍以体积形式存储。

\subsection{控制体积分}

对式\eqref{eq:governing}在 $\Omega_P$ 上积分：

\begin{equation}
  \int_{\Omega_P}
  \frac{\partial\phi}{\partial t}\dd V
  +
  \int_{\Omega_P}
  \nabla\cdot(\bm u\phi)\dd V
  =0.
\end{equation}

固定网格下可交换时间导数与积分：

\begin{equation}
  \int_{\Omega_P}
  \frac{\partial\phi}{\partial t}\dd V
  =
  \frac{\dd}{\dd t}
  \int_{\Omega_P}\phi\dd V
  =
  V_P\frac{\dd\phi_P}{\dd t}.
\end{equation}

利用高斯散度定理：

\begin{equation}
  \int_{\Omega_P}
  \nabla\cdot(\bm u\phi)\dd V
  =
  \int_{\partial\Omega_P}
  \bm u\phi\cdot\bm n_P\dd S.
\end{equation}

将单元边界拆为面集合 $\mathcal F(P)$，采用面中心积分：

\begin{equation}
  \int_{\partial\Omega_P}
  \bm u\phi\cdot\bm n_P\dd S
  \approx
  \sum_{f\in\mathcal F(P)}
  (\bm u_f\cdot\Sf_{P,f})\phi_f.
\end{equation}

定义相对于单元 $P$ 的有向面通量

\begin{equation}
  F_{P,f}
  =
  \bm u_f\cdot\Sf_{P,f},
\end{equation}

得到半离散形式

\begin{equation}
  \boxed{
  V_P\frac{\dd\phi_P}{\dd t}
  +
  \sum_{f\in\mathcal F(P)}
  F_{P,f}\phi_f
  =0
  }.
  \label{eq:semidiscrete}
\end{equation}

等价地，

\begin{equation}
  \boxed{
  \frac{\dd\phi_P}{\dd t}
  =
  -\frac1{V_P}
  \sum_{f\in\mathcal F(P)}
  F_{P,f}\phi_f
  }.
\end{equation}

它称为半离散形式，因为空间已经离散为有限个单元和面，而时间仍连续。整个问题成为ODE系统

\begin{equation}
  \frac{\dd\bm\phi}{\dd t}
  =
  \bm L_h(\bm\phi).
\end{equation}

\subsection{内部面守恒}

设内部面 $f$ 的owner为 $P$，neighbour为 $N$，有

\begin{equation}
  \Sf_{N,f}=-\Sf_{P,f},
\end{equation}

所以

\begin{equation}
  F_{N,f}=-F_{P,f}.
\end{equation}

只要两个单元使用同一个数值面值 $\phi_f$，该面在两个控制体方程中的贡献正好相反：

\begin{equation}
  F_{P,f}\phi_f+F_{N,f}\phi_f=0.
\end{equation}

对所有单元求和时，全部内部面通量相消，只剩边界通量。这就是有限体积法全局守恒的离散来源。若程序中总质量持续漂移，应优先检查共享面是否只计算一次并以相反符号加入owner与neighbour。

\section{空间离散：面值重构}

\subsection{为什么需要面值}

未知量 $\phi_P$ 是单元平均，但输运发生在面上，因此式\eqref{eq:semidiscrete}需要 $\phi_f$。不同空间格式的本质就是如何从周围单元值重构 $\phi_f$。

\subsection{一阶迎风}

对相对于owner单元的面通量

\begin{equation}
  F_f=\bm u_f\cdot\Sf_f,
\end{equation}

一阶迎风定义

\begin{equation}
  \phi_f^{\mathrm{up}}
  =
  \begin{cases}
  \phi_P,&F_f\geq0,\\
  \phi_N,&F_f<0.
  \end{cases}
  \label{eq:first_upwind}
\end{equation}

$F_f\geq0$ 表示信息从owner流向neighbour，所以面值取owner；$F_f<0$ 表示信息从neighbour流向owner，所以面值取neighbour。

将式\eqref{eq:first_upwind}代入半离散方程：

\begin{equation}
  \frac{\dd\phi_P}{\dd t}
  =
  -\frac1{V_P}
  \sum_fF_f\phi_f^{\mathrm{up}}.
  \label{eq:upwind_semidiscrete}
\end{equation}

一阶迎风把单元内的解视为分片常数，光滑区空间截断误差为 $O(h)$。其优势是简单、守恒、稳定性较好，并且在适当CFL下具有离散有界性；主要缺点是人工耗散明显。

\subsection{二阶迎风与线性重构}

设上游单元为 $U$。二阶线性迎风可写成

\begin{equation}
  \phi_f^{(2)}
  =
  \phi_U
  +
  \nabla\phi_U\cdot(\xf_f-\xf_U).
  \label{eq:linear_upwind}
\end{equation}

在光滑网格和光滑解上，如果梯度具有足够精度，则面值重构误差可达到 $O(h^2)$。与一阶迎风相比，它对正弦波、圆锥和光滑凸峰的幅值保持通常更好。

但是线性外推不保证

\begin{equation}
  \min(\phi_U,\phi_D)
  \leq
  \phi_f^{(2)}
  \leq
  \max(\phi_U,\phi_D),
\end{equation}

因此在间断附近可能产生过冲、欠冲、负值和波纹。

\subsection{限制器}

可使用限制系数 $\psi_f\in[0,1]$：

\begin{equation}
  \phi_f
  =
  \phi_U
  +
  \psi_f
  \nabla\phi_U\cdot(\xf_f-\xf_U).
\end{equation}

\begin{itemize}
  \item 光滑区域中 $\psi_f\approx1$，接近二阶重构；
  \item 间断或局部极值附近 $\psi_f$减小；
  \item 必要时 $\psi_f\approx0$，局部退化为一阶迎风。
\end{itemize}

所以“二阶格式”对开槽圆盘不可能在所有位置维持二阶。限制器主动降阶是为了避免非物理振荡，而不是程序失败。

\section{显式时间推进与Courant条件}

\subsection{向前欧拉}

对半离散ODE使用向前欧拉：

\begin{equation}
  \frac{\phi_P^{n+1}-\phi_P^n}{\Delta t}
  =
  -\frac1{V_P}
  \sum_fF_f^n\phi_f^n.
\end{equation}

因此

\begin{equation}
  \boxed{
  \phi_P^{n+1}
  =
  \phi_P^n
  -
  \frac{\Delta t}{V_P}
  \sum_fF_f^n\phi_f^n
  }.
  \label{eq:forward_euler}
\end{equation}

右端全部使用第 $n$ 层已知量，所以不需要求解全局代数方程组。

\subsection{一阶迎风下的凸组合}

对不可压缩速度场，单元面通量满足

\begin{equation}
  \sum_fF_f=0.
\end{equation}

将出流面与入流面分开，一阶迎风更新可写为

\begin{equation}
\begin{aligned}
  \phi_P^{n+1}
  =&
  \left(
  1-\frac{\Delta t}{V_P}
  \sum_{F_f>0}F_f
  \right)\phi_P^n\\
  &+
  \frac{\Delta t}{V_P}
  \sum_{F_f<0}(-F_f)\phi_N^n.
\end{aligned}
\label{eq:convex_update}
\end{equation}

若

\begin{equation}
  \frac{\Delta t}{V_P}
  \sum_{F_f>0}F_f
  \leq1,
  \label{eq:local_cfl}
\end{equation}

则式\eqref{eq:convex_update}中的系数非负且总和为 $1$，新值是上游值的凸组合。因此如果初值和入流边界位于区间 $[m,M]$，数值解也不会轻易越出该区间。这解释了一阶迎风的有界性。

\subsection{OpenFOAM式Courant数}

常用局部定义为

\begin{equation}
  \Co_P
  =
  \frac{\Delta t}{2V_P}
  \sum_f|F_f|.
\end{equation}

不可压缩条件下，入流绝对通量之和等于出流通量之和，因此

\begin{equation}
  \frac12\sum_f|F_f|
  =
  \sum_{F_f>0}F_f.
\end{equation}

所以

\begin{equation}
  \Co_P
  =
  \frac{\Delta t}{V_P}
  \sum_{F_f>0}F_f.
\end{equation}

最大Courant数为

\begin{equation}
  \Cmax=\max_P\Co_P.
\end{equation}

题目规定 $\Cmax=0.2$，意味着应根据最小单元、局部速度和面通量自动调整 $\Delta t$。

\subsection{规则方格上的具体关系}

对 $\bm u=(1,1)$、$\Delta x=\Delta y=h$ 的规则方格，

\begin{equation}
  \Co
  =
  \frac{\Delta t}{h}
  +
  \frac{\Delta t}{h}
  =
  \frac{2\Delta t}{h}.
\end{equation}

因此

\begin{equation}
  \Delta t
  =
  \frac{\Co h}{2}.
\end{equation}

$\Co=0.2$ 时，

\begin{equation}
  \Delta t=0.1h.
\end{equation}

三角形或非均匀网格中不能仅用统一 $h$ 估算，必须根据每个单元的 $V_P$ 与面通量计算局部 $\Co_P$。

\subsection{二阶空间格式的时间推进}

向前欧拉是一阶时间格式。即使空间重构是二阶，在固定CFL下通常有 $\Delta t=O(h)$，总体误差仍可能被 $O(\Delta t)$ 限制为一阶。

此外，纯二阶迎风空间算子与向前欧拉的稳定性并不像一阶迎风那样自然。为了可靠比较二阶空间格式，通常更合适使用显式RK2或SSP-RK3。若必须使用向前欧拉，则需要进行独立的稳定性试验，并采用足够小的时间步，不能直接假定与一阶迎风相同的CFL上限。

\section{数值耗散与数值色散}

\subsection{一阶迎风修正方程的推导}

考虑一维正速度 $a>0$ 的对流方程

\begin{equation}
  \phi_t+a\phi_x=0.
\end{equation}

一阶迎风加向前欧拉为

\begin{equation}
  \frac{\phi_i^{n+1}-\phi_i^n}{\Delta t}
  +
  a\frac{\phi_i^n-\phi_{i-1}^n}{\Delta x}
  =0.
  \label{eq:ftbs}
\end{equation}

在 $(x_i,t_n)$ 处展开：

\begin{equation}
  \frac{\phi_i^{n+1}-\phi_i^n}{\Delta t}
  =
  \phi_t
  +\frac{\Delta t}{2}\phi_{tt}
  +O(\Delta t^2),
\end{equation}

\begin{equation}
  \frac{\phi_i^n-\phi_{i-1}^n}{\Delta x}
  =
  \phi_x
  -\frac{\Delta x}{2}\phi_{xx}
  +O(\Delta x^2).
\end{equation}

代入式\eqref{eq:ftbs}：

\begin{equation}
  \phi_t+a\phi_x
  =
  \frac{a\Delta x}{2}\phi_{xx}
  -
  \frac{\Delta t}{2}\phi_{tt}
  +O(\Delta x^2,\Delta t^2).
\end{equation}

由原方程 $\phi_t=-a\phi_x$，再次对时间求导：

\begin{equation}
  \phi_{tt}=a^2\phi_{xx}.
\end{equation}

所以

\begin{equation}
  \boxed{
  \phi_t+a\phi_x
  =
  \frac{a\Delta x}{2}
  \left(1-\frac{a\Delta t}{\Delta x}\right)
  \phi_{xx}
  +\cdots
  }.
\end{equation}

令

\begin{equation}
  C=\frac{a\Delta t}{\Delta x},
\end{equation}

人工扩散系数为

\begin{equation}
  \nu_{\mathrm{num}}
  =
  \frac{a\Delta x}{2}(1-C).
  \label{eq:numerical_diffusion}
\end{equation}

这解释了为什么一阶迎风把纯对流问题算成“对流加人工扩散”。

\subsection{耗散在图上的对应}

\begin{itemize}
  \item 正弦波振幅降低：人工扩散优先削弱高频分量；
  \item 开槽圆盘边缘模糊：间断被扩散成有限宽度过渡层；
  \item 槽口填平：槽口两侧的高值通过人工扩散进入低值区域；
  \item 圆锥峰值降低、底部变宽：质量被摊到更多单元；
  \item 光滑凸峰变矮、变宽：曲率较大的区域受到更明显的扩散。
\end{itemize}

如果总质量守恒而最大值下降，说明标量没有消失，而是被重新分布。

\subsection{色散与相位误差}

数值色散主要表现为不同波数的传播速度不正确。图像上的典型表现包括：

\begin{itemize}
  \item 正弦波波峰整体偏移；
  \item 复杂轮廓旋转角度稍有超前或滞后；
  \item 间断前后出现正负交替波纹；
  \item 高阶无限制格式产生过冲和欠冲。
\end{itemize}

耗散主要改变幅值，色散主要改变相位。实际结果中两者可能同时存在。

\subsection{Courant数与耗散的关系}

式\eqref{eq:numerical_diffusion}说明，在一维一阶迎风加向前欧拉、固定网格的条件下，减小 $C$ 会使 $\nu_{\mathrm{num}}$趋向 $a\Delta x/2$。所以“时间步越小，最终轮廓一定越清晰”并不成立：时间误差减小了，但空间迎风耗散仍然存在，并且固定物理时间内需要执行更多次迎风更新。

增大 $C$ 在稳定范围内可能减小部分耗散，但会降低稳定裕度，并增加时间截断误差。该结论是特定格式的修正方程结论，不能不加分析地推广到所有二阶格式和非结构网格。

\section{网格分辨率\texorpdfstring{$N$}{N}及计算成本}

\subsection{\texorpdfstring{$N$}{N}表示区间数}

若 $N$ 是每个方向的区间数，单位正方形上

\begin{equation}
  h=\frac1N.
\end{equation}

规则四边形单元数为

\begin{equation}
  N_{\mathrm{quad}}=N^2.
\end{equation}

若将每个四边形拆成两个三角形，则

\begin{equation}
  N_{\mathrm{tri}}=2N^2.
\end{equation}

\subsection{\texorpdfstring{$N$}{N}表示顶点数}

若题目说每条边有 $N$ 个顶点，则区间数为 $N-1$：

\begin{equation}
  N_{\mathrm{quad}}=(N-1)^2,
  \qquad
  N_{\mathrm{tri}}=2(N-1)^2.
\end{equation}

所以“边上100个顶点”严格对应99个区间。

\subsection{增大\texorpdfstring{$N$}{N}的影响}

\begin{longtable}{L{0.37\linewidth}L{0.51\linewidth}}
\toprule
变化 & 原因与结果\\
\midrule
\endfirsthead
\toprule
变化 & 原因与结果\\
\midrule
\endhead
$h$减小 & 网格空间分辨率提高\\
一阶迎风耗散减小 & 人工扩散的主导量通常与 $h$ 成正比\\
光滑区二阶空间误差更快下降 & 理想情况下为 $O(h^2)$\\
开槽边缘更清晰 & 过渡层占据的物理宽度减小\\
圆锥与凸峰峰值保持更好 & 人工扩散减弱\\
允许时间步减小 & 固定CFL下 $\Delta t=O(h)$\\
内存增加 & 单元与面数量约为 $O(N^2)$\\
时间步数增加 & 固定终止时间下约为 $O(N)$\\
总工作量增加 & 显式二维计算约为 $O(N^3)$\\
\bottomrule
\end{longtable}

$N$加倍时，二维单元数约增至4倍，固定CFL下时间步数约增至2倍，所以总工作量可能约增至8倍。

\section{三角形与四边形网格}

\subsection{几何与拓扑差异}

四边形规则网格的面方向和邻接结构较简单；三角形网格更适合复杂几何，但在相同背景节点间距下通常拥有更多单元和面。

数值通量由

\begin{equation}
  F_f=\bm u_f\cdot\Sf_f
\end{equation}

决定。两类网格的面方向、面长度、单元中心位置、上游邻居与局部CFL均不同，所以误差不仅由单一标量 $h$ 决定，还具有方向性。

\subsection{方向性耗散}

在多维和非结构网格上，人工扩散更接近与网格和速度方向相关的扩散张量。图像中可能表现为：

\begin{itemize}
  \item 某些方向比其他方向更模糊；
  \item 旋转后的圆盘略呈椭圆；
  \item 槽口两侧模糊程度不对称；
  \item 三角形与四边形误差阶相近，但误差常数不同。
\end{itemize}

不能预先断定三角形一定更差。公平比较至少应说明采用的是：

\begin{itemize}
  \item 相同背景节点间距；
  \item 相近特征网格尺寸 $h$；
  \item 或相近自由度与计算成本。
\end{itemize}

这些比较标准不能同时完全相同，因此报告中必须明确选择。

\section{误差、收敛阶与诊断量}

\subsection{单元平均与解析比较值}

有限体积未知量是单元平均

\begin{equation}
  \phi_P
  =
  \frac1{V_P}\int_{\Omega_P}\phi\dd V.
\end{equation}

严格比较时，解析值也应取单元平均

\begin{equation}
  \phi_{P,\mathrm{exact}}
  =
  \frac1{V_P}
  \int_{\Omega_P}
  \phi_{\mathrm{exact}}\dd V.
\end{equation}

实际程序常用单元中心解析值近似该平均值。对于一阶格式通常足够，但在验证高阶收敛时，低阶解析积分可能污染误差阶，应采用高阶积分或解析单元平均。

\subsection{误差范数}

设数值解与解析解分别为 $\phi_P^{\mathrm{num}}$ 和 $\phi_P^{\mathrm{exact}}$。体积加权相对误差可定义为

\begin{equation}
  E_{L_1}
  =
  \frac{
  \sum_P
  |\phi_P^{\mathrm{num}}-\phi_P^{\mathrm{exact}}|V_P
  }{
  \sum_P
  |\phi_P^{\mathrm{exact}}|V_P
  },
\end{equation}

\begin{equation}
  E_{L_2}
  =
  \left[
  \frac{
  \sum_P
  (\phi_P^{\mathrm{num}}-\phi_P^{\mathrm{exact}})^2V_P
  }{
  \sum_P
  (\phi_P^{\mathrm{exact}})^2V_P
  }
  \right]^{1/2},
\end{equation}

\begin{equation}
  E_{L_\infty}
  =
  \frac{
  \max_P|\phi_P^{\mathrm{num}}-\phi_P^{\mathrm{exact}}|
  }{
  \max_P|\phi_P^{\mathrm{exact}}|
  }.
\end{equation}

$L_1$反映总体平均偏差，$L_2$更强调较大误差，$L_\infty$反映最坏单元误差。

\subsection{收敛阶}

两套网格尺寸为 $h_1>h_2$，误差为 $E_1,E_2$，观测阶为

\begin{equation}
  p
  =
  \frac{\log(E_1/E_2)}
  {\log(h_1/h_2)}.
\end{equation}

若每个方向加密2倍，则 $h_2=h_1/2$：

\begin{equation}
  p
  =
  \frac{\log(E_1/E_2)}{\log2}.
\end{equation}

二维中若只知道总单元数 $N_c$，通常

\begin{equation}
  h\sim N_c^{-1/2}.
\end{equation}

因此不能直接把总单元数之比当作网格尺寸之比。

\subsection{总体守恒}

定义离散总量

\begin{equation}
  M(t)=\sum_PV_P\phi_P(t).
\end{equation}

周期边界或无净边界通量时，守恒格式应使 $M(t)$保持到舍入误差范围。

正弦波的有符号总量理论上为零，因此相对质量误差的分母可能接近零。可以报告

\begin{equation}
  \frac{|M(t)-M(0)|}
  {\sum_PV_P|\phi_P(0)|}
\end{equation}

或直接报告均值漂移。复杂轮廓初值非负，适合使用普通相对质量误差。

\subsection{有界性与极值}

记录

\begin{equation}
  \phi_{\max}(t)=\max_P\phi_P(t),
  \qquad
  \phi_{\min}(t)=\min_P\phi_P(t).
\end{equation}

对非负复杂轮廓，一阶迎风和合适CFL应大致保持非负。如果出现显著负值或超过初始最大值，应检查高阶重构、限制器、CFL和边界处理。

\subsection{正弦波的幅值和相位}

仅看 $L_1$无法区分耗散与相位误差。可沿一条截线叠加数值解与精确解，也可将数值场投影到对应正弦和余弦模态。若数值波峰位置正确但振幅下降，主要是耗散；若振幅基本保持但波峰横向偏移，主要是色散或时间相位误差。

\section{实验设计与理论预期}

\subsection{正弦波网格收敛实验}

建议固定：

\begin{itemize}
  \item 空间格式：一阶迎风；
  \item 时间格式：向前欧拉；
  \item $\Cmax=0.2$；
  \item 终止时间：$t=1$；
  \item 周期边界。
\end{itemize}

改变每方向分辨率，例如

\begin{equation}
  N=20,\ 40,\ 80,\ 160,
\end{equation}

分别构造四边形和三角形网格。每组实验记录：

\begin{equation}
  N_c,\ h,\ \Delta t,\ E_{L_1},\ E_{L_2},\
  E_{L_\infty},\ p,\ \phi_{\max},\ \phi_{\min}.
\end{equation}

理论预期是：网格加密后误差下降，一阶迎风加向前欧拉、固定CFL时，光滑正弦波的总体观测阶应趋近1。粗网格上振幅衰减更明显。

\subsection{Courant敏感性实验}

固定网格与空间格式，可比较

\begin{equation}
  \Cmax=0.05,\ 0.1,\ 0.2,\ 0.4
\end{equation}

等位于稳定范围内的值。需要同时观察误差、振幅、步数和运行成本。不能仅凭“时间步更小”判断结果一定更好；应结合时间误差与空间迎风耗散解释。

\subsection{一阶与二阶空间格式}

为了公平比较：

\begin{itemize}
  \item 使用同一网格与同一初边值条件；
  \item 一阶迎风可配向前欧拉作为基准；
  \item 二阶线性迎风更适合配显式RK2或SSP-RK3；
  \item 对复杂轮廓应同时测试无限制与带限制器版本；
  \item 报告误差、最大最小值、质量和图像，而不能只比较视觉清晰度。
\end{itemize}

理论预期是：二阶格式在正弦波和光滑凸峰上耗散更小；无限制二阶格式在开槽圆盘附近可能产生过冲和波纹；带限制器格式在清晰度与有界性之间折中。

\subsection{复杂轮廓正式实验}

按题目要求，计算区域每条边布置100个顶点，分别构造三角形与四边形网格，计算到

\begin{equation}
  t=2\pi.
\end{equation}

至少输出初始轮廓、最终轮廓、二者叠加图、差值图、质量曲线、最大最小值曲线。若进行扩展实验，可比较不同 $N$、不同格式和不同CFL，但不应因此延迟基准算例的完成。

\section{结果图的读取与因果诊断}

\begin{longtable}{L{0.38\linewidth}L{0.50\linewidth}}
\toprule
观察到的现象 & 优先解释与检查方向\\
\midrule
\endfirsthead
\toprule
观察到的现象 & 优先解释与检查方向\\
\midrule
\endhead
正弦波振幅降低 & 一阶迎风人工耗散；检查网格分辨率和空间格式\\
正弦波波峰整体偏移 & 相位误差；检查时间格式、空间色散和终止时刻\\
波峰前后出现正负误差条带 & 相位偏移或高阶色散\\
开槽圆盘边缘模糊 & 间断被人工扩散成有限宽度过渡层\\
槽口变窄或填平 & 槽口两侧高值向低值区数值扩散\\
圆锥峰值下降且底部变宽 & 质量守恒条件下的人工扩散重分布\\
光滑凸峰变矮、变宽 & 光滑结构上的数值耗散\\
间断附近出现波纹 & 高阶无限制重构的过冲、欠冲与色散\\
出现负值或超过初始最大值 & 有界性破坏；检查限制器和CFL\\
所有数值迅速增大 & 显式时间推进失稳或符号错误\\
轮廓整体转角不等于 $2\pi$ & 速度场、时间推进、终止时间或输出时间错误\\
旋转半径改变 & 面速度、通量插值或几何装配错误\\
圆形变成椭圆或方向性模糊 & 网格各向异性与方向性人工扩散\\
总质量持续变化 & 内部面未成对装配或边界通量处理错误\\
误差集中在周期边界 & 周期面配对、方向或面值传递错误\\
网格加密但误差不下降 & 时间误差、边界误差、解析比较值或代码错误主导\\
误差下降后进入平台 & 可能达到时间误差、边界误差或舍入误差主导区\\
\bottomrule
\end{longtable}

不同实验的对比图必须使用相同色标、等值线级别、坐标范围、纵横比例和输出时刻。自动缩放色标可能掩盖振幅衰减、过冲和负值。

\section{完整求解器算法}

\subsection{几何预处理}

对每个单元和面准备：

\begin{itemize}
  \item 单元体积或面积 $V_P$；
  \item 面面积向量 $\Sf_f$；
  \item owner与neighbour；
  \item 单元中心 $\xf_P$ 与面中心 $\xf_f$；
  \item 边界面类型与周期配对。
\end{itemize}

\subsection{时间循环}

\begin{lstlisting}[style=algorithm]
read mesh and boundary information
construct cell volumes and face area vectors
initialise phi from the selected test case
set time t = 0

while t < endTime:
    evaluate velocity at cells or faces
    compute face flux F_f = u_f dot S_f
    compute local Courant numbers
    choose dt so that max(Co_P) <= CoMax
    shorten the last dt if t + dt > endTime

    reconstruct phi_f on every internal and boundary face
    assemble the conservative residual:
        owner     receives -F_f * phi_f
        neighbour receives +F_f * phi_f

    update every cell explicitly:
        phi_P <- phi_P + dt * residual_P / V_P

    apply or correct boundary conditions
    update time
    write fields and diagnostics at selected times

compute exact solution, errors and convergence data
generate contour, line, error and diagnostic plots
\end{lstlisting}

\subsection{OpenFOAM中的对应概念}

\begin{longtable}{L{0.34\linewidth}L{0.54\linewidth}}
\toprule
数学对象 & OpenFOAM中的对应\\
\midrule
\endfirsthead
\toprule
数学对象 & OpenFOAM中的对应\\
\midrule
\endhead
单元标量 $\phi_P$ & \texttt{volScalarField}\\
面通量 $F_f$ & \texttt{surfaceScalarField}\\
单元体积 $V_P$ & \texttt{mesh.V()}\\
面面积向量 $\Sf_f$ & \texttt{mesh.Sf()}\\
owner与neighbour & \texttt{faceOwner()}与\texttt{faceNeighbour()}\\
显式散度 & \texttt{fvc::div(flux, phi)}\\
一阶迎风 & \texttt{Gauss upwind}\\
线性迎风 & \texttt{Gauss linearUpwind grad(phi)}\\
周期边界 & \texttt{cyclic} patch\\
最大Courant数 & \texttt{maxCo}与时间步控制\\
\bottomrule
\end{longtable}

若使用\texttt{fvc::div}，空间算子是显式返回的单元场；若使用\texttt{fvm::div}，则通常构造隐式矩阵。第一题要求全显式基准时，应明确程序实际采用了哪一种。

\section{结果交付结构}

一个能够证明理解与实现能力的第一题交付建议包含：

\begin{enumerate}
  \item 自包含题目与符号说明；
  \item 连续方程、特征线和两个精确运动的说明；
  \item 有限体积半离散推导；
  \item 一阶迎风与显式时间推进；
  \item CFL和有界性分析；
  \item 求解器代码与编译运行说明；
  \item 四边形和三角形网格说明；
  \item 正弦波网格收敛表和误差图；
  \item 复杂轮廓初始、最终、叠加和差值图；
  \item 质量、最大值、最小值诊断；
  \item 数值耗散、色散和网格方向性解释；
  \item 已知限制、题面歧义和下一步改进。
\end{enumerate}

实际数值表格中必须填写程序运行结果，不能用理论阶数代替观测阶，也不能根据预期趋势虚构误差。

\section{老师通过第一题可能考察什么}

该题的结构表明，考察重点不是单一公式，而是以下能力：

\begin{itemize}
  \item 能否将守恒PDE转化为有限体积离散；
  \item 能否理解面通量、上游方向与边界信息传播；
  \item 能否从离散式实现可运行求解器；
  \item 能否用解析解、误差和收敛阶验证程序；
  \item 能否用复杂轮廓识别耗散、振荡和有界性问题；
  \item 能否区分网格、时间步和空间格式的影响；
  \item 能否面对不完整条件时明确假设，而不是隐式猜测；
  \item 能否把Agent或VibeFlow生成的配置、代码和结果进行独立审查。
\end{itemize}

完成第一题的真正标准不是“程序输出了一张图”，而是能够回答：

\begin{equation}
\boxed{
\text{为什么使用这个离散，为什么得到这种图，
误差来自哪里，怎样证明程序正在收敛。}
}
\end{equation}

\section{结论}

第一题构成了一个完整的数值求解器验证闭环。连续方程告诉我们精确运动不应发生扩散；有限体积法把守恒律转化为单元面通量平衡；一阶与二阶重构决定空间精度、耗散、有界性和振荡；显式时间格式与Courant数决定时间误差和稳定性；$N$与网格拓扑决定分辨率、方向性误差和计算成本；两个算例分别提供光滑定量验证与复杂结构定性诊断。

真正掌握该题，意味着能够把

\begin{equation}
\boxed{
\text{方程}
\rightarrow
\text{离散}
\rightarrow
\text{代码}
\rightarrow
\text{实验}
\rightarrow
\text{误差}
\rightarrow
\text{原因}
}
\end{equation}

连成一条可复现、可解释、可检查的完整链路。

\end{document}
